\documentclass[10pt]{article}
\usepackage[T1]{fontenc}
\usepackage[utf8]{inputenc}
\usepackage[french]{babel}
\usepackage{amsmath,amssymb,mathrsfs,mathpazo}
\usepackage{geometry}
\geometry{a4paper,margin=2cm}
\usepackage{graphicx,array,tabularx,tabularray,colortbl,multirow,multicol,booktabs,enumitem,pifont,ulem,fancybox,diagbox,cancel}
\usepackage{pstricks,pstricks-add,pst-plot,pst-tree,pst-node,pst-text,pst-func,pst-math,pst-bspline,pst-3dplot,pst-eucl,pst-solides3d}
\let\pscalculate\undefined
\usepackage{pst-bezier}
\usepackage{tikz,pgfplots,tkz-tab}
\usetikzlibrary{arrows,arrows.meta,patterns,calc,positioning,trees,decorations.markings,shapes.geometric,shapes.misc,shapes.arrows,intersections,decorations.pathmorphing,decorations.pathreplacing,backgrounds,shadows,babel}
\pgfplotsset{compat=1.18}
\usepgfplotslibrary{fillbetween}
\usepackage[np]{numprint}
\usepackage{listings,makecell,pas-tableur,twoopt}
\let\barre\undefined
\usepackage[pstricks]{tablvar}
\definecolor{darkgreen}{RGB}{0,100,0}
\providecommand{\red}{\color{red}}
\providecommand{\blue}{\color{blue}}
\providecommand{\green}{\color{green}}
\providecommand{\black}{\color{black}}
\providecommand{\white}{\color{white}}
\providecommand{\gray}{\color{gray}}
\providecommand{\cyan}{\color{cyan}}
\providecommand{\magenta}{\color{magenta}}
\providecommand{\yellow}{\color{yellow}}
\providecommand{\violet}{\color{violet}}
\providecommand{\orange}{\color{orange}}
\providecommand{\euro}{\texteuro}
\providecommand{\EUR}{\texteuro}
\providecommand{\ieme}{\textsuperscript{e}}
\providecommand{\degres}{\ensuremath{^\circ}}
\providecommand{\degree}{\ensuremath{^\circ}}
\providecommand{\vv}[1]{\overrightarrow{#1}}
\providecommand{\Coord}[1]{#1}
\providecommand{\decosix}{}
\providecommand{\starredbullet}{\ensuremath{\bullet}}
\pagestyle{empty}
\setlength{\parindent}{0pt}
\usepackage{adjustbox}\DeclareUnicodeCharacter{2236}{:}\DeclareUnicodeCharacter{2190}{\ensuremath{\leftarrow}}
\providecommand{\R}{\mathbb{R}}
\providecommand{\N}{\mathbb{N}}
\providecommand{\Z}{\mathbb{Z}}
\providecommand{\Q}{\mathbb{Q}}
\providecommand{\D}{\mathbb{D}}
\providecommand{\C}{\mathbb{C}}
\providecommand{\e}{{\mathrm{e}}}
\providecommand{\vect}[1]{\overrightarrow{#1}}
\providecommand{\vectt}[1]{\overrightarrow{\mathrm{#1}}}
\providecommand{\vvt}[1]{\overrightarrow{\mathrm{#1}}}
\providecommand{\np}[2][]{{#2~#1}}
\providecommand{\vv}[1]{{\overrightarrow{#1}}}
\renewcommand{\d}{{\mathrm{d}}}
\providecommand{\degres}{{^\circ}}
\providecommand{\degre}{{^\circ}}
\providecommand{\blue}{\color{blue}}
\providecommand{\red}{\color{red}}
\providecommand{\enskip}{\quad}
\providecommand{\llbracket}{\lbrack\!\lbrack}
\providecommand{\rrbracket}{\rbrack\!\rbrack}
\providecommand{\euro}{{\text{€}}}
\providecommand{\no}{{\text{n°}}}
\providecommand{\interval}[2]{{[#1;#2]}}
\providecommand{\intervFO}[2]{{[#1;#2[}}
\providecommand{\dotfill}{\ldots}

\providecommand{\ds}{\displaystyle}
\providecommand{\textding}[1]{#1}
\providecommand{\Oij}{\(\mathrm{O};\vec{\imath},\vec{\jmath}\)}
\providecommand{\Oijk}{\(\mathrm{O};\vec{\imath},\vec{\jmath},\vec{k}\)}
\providecommand{\Ouv}{\(\mathrm{O};\vec{u},\vec{v}\)}
\providecommand{\intervFF}[2]{\ensuremath{[#1;#2]}}
\renewcommand{\euro}{\ifmmode\text{\texteuro}\else\texteuro\fi}
\DeclareUnicodeCharacter{00B0}{\ifmmode^{\circ}\else\textdegree\fi}
\DeclareUnicodeCharacter{20AC}{\ifmmode\text{\texteuro}\else\texteuro\fi}
% Énoncés : Éducation nationale. Transcription et corrigés : APMEP, Denis Vergès et contributeurs. Découpage et classement : CoopMaths, Sébastien Lozano, Éric Elter et contributeurs. Adaptation de la présentation : Maths Explorer.
% CC BY-SA : https://coopmaths.fr/www/about/
% https://coopmaths.fr/alea/static/e3c/2020/tex/e3c_2020_00_sujet12_3.tex
\begin{document}
\noindent\textit{APMEP / CoopMaths. CC BY-SA. Présentation adaptée par Maths Explorer.}\par\medskip
\medskip
Dans la figure ci-dessous, on a tracé $\mathcal{C}_f$ la courbe représentative d'une fonction $f$  définie et dérivable sur $\R$ ainsi que les tangentes à $\mathcal{C}_f$ aux points d'abscisses respectives $- 2$,\:$- 1$ et $0$.

\begin{center} 

\begin{tikzpicture}
  \begin{axis}[
    width=15cm, 
    height=10cm, 
    grid=both,
    grid style={line width=.1pt, draw=gray!50}, 
    minor grid style={line width=.05pt, draw=gray!20}, 
    minor tick num=4, 
    axis equal image, 
    axis lines=middle,
    xlabel={$x$}, ylabel={$y$},
    xmin=-3.5, xmax=3.5, 
    ymin=-3, ymax=7, 
    xtick={-3,-2,...,3}, 
    ytick={-3,-2,...,7}, 
    enlargelimits=false,
    domain=-3:1,
    samples=200
  ]
    \addplot[red, thick, no markers] {x^3 + 3*x^2 + 2*x + 1};
    \addplot[blue, thick, domain=-3.2:-0.8, <->] {
      (-2)^3 + 3*(-2)^2 + 2*(-2) + 1 + (3*(-2)^2 + 6*(-2) + 2)*(x - -2)
    };
    \addplot[blue, thick, domain=-2.2:0.2, <->] {
      (-1)^3 + 3*(-1)^2 + 2*(-1) + 1 + (3*(-1)^2 + 6*(-1) + 2)*(x - -1)
    };
    \addplot[ blue, thick, domain=-1.2:1.2, <->] {
      (0)^3 + 3*(0)^2 + 2*(0) + 1 + (3*(0)^2 + 6*(0) + 2)*(x - 0)
    };
  \end{axis}
\end{tikzpicture}

\end{center}


\begin{enumerate}
\item Recopier sur la copie en le complétant le tableau de valeurs ci-dessous. 

\begin{center}
\psset{unit=1cm,arrowsize=2pt 4}
\begin{tabularx}{0.5\linewidth}{|*{3}{>{\centering \arraybackslash}X|}}\hline
$x$		&$-1$	& $0$\\ \hline
$f(x)$	&		&\\ \hline
$f'(x)$	&		&\\ \hline
\end{tabularx}
\end{center}

On admet que la fonction $f$ est définie sur $\R$ par:

\[f(x) = x^3 +3x^2 +2x +1.\]

\item 
	\begin{enumerate}
		\item Calculer $f'(x)$, pour tout réel $x$.
		\item Résoudre dans $\R$ l'équation: $f'(x) = 0$.
	\end{enumerate}
\item Dresser le tableau de variations de la fonction $f$.

\item Le point S$(-4~;~-3)$ appartient-il à la tangente à la courbe représentative de $f$ au point d'abscisse $x = -2$ ?
\end{enumerate}

\bigskip

\end{document}
